% Motor: XeLaTeX, fijado en .latexmkrc y en la receta del editor.
\documentclass[aspectratio=169,10pt,c]{beamer}

% !TEX root = slides.tex
% Preambulo comun: estilo UTP, colores, fuentes, plantillas de beamer y
% comandos de apoyo didactico. Se carga desde slides.tex tras \documentclass.

% La opción "c" de beamer reparte el espacio libre 1:1.5 (arriba:abajo), por lo
% que el contenido queda por encima del centro. Igualamos ambos rellenos.
\makeatletter
\define@key{beamerframe}{c}[true]{%
  \beamer@frametopskip=0pt plus 1fill\relax%
  \beamer@framebottomskip=0pt plus 1fill\relax%
  \beamer@frametopskipautobreak=0pt plus .5\paperheight\relax%
  \beamer@framebottomskipautobreak=0pt plus .5\paperheight\relax%
}
\makeatother

\usepackage{graphicx}
\usepackage{tikz}
\usepackage{xcolor}
\usepackage{iftex}

\ifPDFTeX
  \usepackage[utf8]{inputenc}
  \usepackage[T1]{fontenc}
  \usepackage{lmodern}          % tipografia escalable para el respaldo con pdflatex
  \renewcommand{\familydefault}{\sfdefault}
\else
  \usepackage{fontspec}
  \IfFontExistsTF{Arial}
    {\setsansfont{Arial}[BoldFont={Arial Bold}]}
    {\setsansfont{Helvetica}[BoldFont={Helvetica Bold}]}
\fi

\usepackage[spanish,es-nodecimaldot]{babel}
\usepackage{microtype}
\usepackage{booktabs}
\usepackage{amssymb}
\usepackage{array}
\usetikzlibrary{arrows.meta,positioning,fit,calc,shapes.geometric}

\graphicspath{{assets/}}

% ---------- Identidad visual UTP ----------
\definecolor{utpRojo}{HTML}{C42B35}
\definecolor{utpNegro}{HTML}{111111}
\definecolor{utpGris}{HTML}{5F5F5F}
\definecolor{utpGrisMedio}{HTML}{9A9A9A}
\definecolor{utpGrisClaro}{HTML}{F1F1F1}
\definecolor{utpNaranja}{HTML}{C1670B}

\usetheme{default}
\setbeamercolor{normal text}{fg=utpNegro,bg=white}
\setbeamercolor{structure}{fg=utpRojo}
\setbeamercolor{frametitle}{fg=utpNegro,bg=white}
\setbeamercolor{title}{fg=utpNegro}
\setbeamercolor{subtitle}{fg=utpNegro}
\setbeamercolor{author}{fg=utpNegro}
\setbeamercolor{date}{fg=utpNegro}
\setbeamercolor{block title}{fg=white,bg=utpNegro}
\setbeamercolor{block body}{fg=utpNegro,bg=utpGrisClaro}
\setbeamercolor{alertblock title}{fg=white,bg=utpRojo}
\setbeamercolor{alertblock body}{fg=utpNegro,bg=utpRojo!7}
\setbeamercolor{exampleblock title}{fg=white,bg=utpGris}
\setbeamercolor{exampleblock body}{fg=utpNegro,bg=utpGris!8}

\setbeamerfont{title}{series=\bfseries,size=\fontsize{19.5}{22.5}\selectfont}
\setbeamerfont{subtitle}{size=\fontsize{11.5}{13.5}\selectfont}
\setbeamerfont{date}{series=\bfseries,size=\fontsize{11.5}{13.5}\selectfont}
\setbeamerfont{author}{size=\fontsize{10.5}{12.5}\selectfont}
\setbeamerfont{frametitle}{series=\bfseries,size=\large}
\setbeamerfont{block title}{series=\bfseries,size=\small}
\setbeamerfont{block body}{size=\normalsize}

\setbeamertemplate{navigation symbols}{}
\setbeamertemplate{headline}{}
\setbeamertemplate{blocks}[default]
\setbeamertemplate{itemize item}{\color{utpRojo}\small\raise.15ex\hbox{$\blacktriangleright$}}
\setbeamertemplate{itemize subitem}{\color{utpGris}\scriptsize$\bullet$}
\setbeamertemplate{enumerate item}{\color{utpRojo}\insertenumlabel.}
\setbeamertemplate{background canvas}{\color{white}\rule{\paperwidth}{\paperheight}}

% Título de diapositiva: negrita negra a la izquierda, logo UTP a la derecha
\setbeamertemplate{frametitle}{%
  \begin{tikzpicture}[remember picture,overlay]
    \node[anchor=north east,inner sep=0pt]
      at ([xshift=-7mm,yshift=-5.5mm]current page.north east)
      {\includegraphics[width=17mm]{utp-logo.png}};
  \end{tikzpicture}%
  \vspace{1.2mm}%
  \begin{minipage}[t]{.80\textwidth}\raggedright\insertframetitle\end{minipage}\par
  \vspace{1.1mm}{\color{utpRojo}\rule{\textwidth}{1.1pt}}
}

% Solo numeración discreta; la plantilla UTP no lleva pie de página
\setbeamertemplate{footline}{%
  \hfill{\color{utpGrisMedio}\scriptsize\insertframenumber}\hspace{6mm}\vspace{2.5mm}%
}

% Portada UTP
\setbeamertemplate{title page}{%
  \begin{tikzpicture}[remember picture,overlay]
    \node[anchor=center,inner sep=0pt] at (current page.center) {%
      \includegraphics[width=\paperwidth,height=\paperheight]{utp-cover-bg.jpeg}%
    };
    \node[anchor=north,align=center,text=utpNegro,inner sep=0pt,
      text width=145mm,font=\usebeamerfont{title}]
      at ([yshift=-25.2mm]current page.north) {\inserttitle};
    \node[anchor=north,align=center,text=utpNegro,inner sep=0pt,
      text width=142mm,font=\usebeamerfont{subtitle}]
      at ([yshift=-38.2mm]current page.north) {\insertsubtitle};
    \node[anchor=north,align=center,text=utpNegro,inner sep=0pt,
      font=\usebeamerfont{date}]
      at ([yshift=-50.4mm]current page.north) {\insertdate};
    \node[anchor=north,align=center,text=utpNegro,inner sep=0pt,
      font=\usebeamerfont{author}]
      at ([yshift=-58.4mm]current page.north) {\insertauthor};
  \end{tikzpicture}%
}

% ---------- Comandos de apoyo didáctico ----------
\newcommand{\concepto}[2]{%
  \begin{block}{#1}\large #2\end{block}}
\newcommand{\pregunta}[1]{%
  \begin{alertblock}{Decide y explica}#1\end{alertblock}}
\newcommand{\analogia}[1]{%
  \begin{exampleblock}{En palabras simples}\small #1\end{exampleblock}}
\newcommand{\enlarealidad}[1]{%
  \begin{block}{Cómo se ve en la realidad}\small #1\end{block}}
\newcommand{\ojo}[1]{%
  \begin{alertblock}{Ojo con esto}\small #1\end{alertblock}}
\newcommand{\apoyo}{{\normalfont\footnotesize\color{utpGris}~· material de apoyo}}
\newcommand{\marca}[1]{\tikz[baseline=(m.base)]{%
  \node[rounded corners=1.5pt,fill=utpGrisClaro,draw=utpGrisMedio,inner xsep=3pt,inner ysep=1.5pt,
        font=\footnotesize] (m) {#1};}}

\tikzset{
  caja/.style={draw=utpNegro, very thick, rounded corners=1.5pt, fill=white,
    align=center, minimum height=0.72cm, inner sep=5pt, font=\small},
  cajaazul/.style={caja, draw=utpGris, fill=white},
  cajaverde/.style={caja, draw=utpNegro, fill=utpGrisClaro},
  cajanaranja/.style={caja, draw=utpRojo!75, fill=utpRojo!7},
  flecha/.style={-{Stealth[length=2.3mm]}, very thick, draw=utpNegro},
  flechavuelta/.style={{Stealth[length=2.3mm]}-{Stealth[length=2.3mm]}, very thick, draw=utpNegro}
}

% ---------- Notación pata de gallo (crow's foot) ----------
% La rotación indica hacia dónde está la entidad: 0 derecha, 180 izquierda,
% 90 arriba, -90 abajo. La pata se abre hacia la entidad del lado "muchos".
\tikzset{
  relacion/.style={thick, draw=utpNegro},
  pics/uno/.style={code={%
    \draw[thick,utpNegro,line cap=round] (0,-1.5mm) -- (0,1.5mm);}},
  pics/muchos/.style={code={%
    \draw[thick,utpNegro,line cap=round,line join=round]
      (2.6mm,1.5mm) -- (0,0) -- (2.6mm,-1.5mm);
    \draw[thick,utpNegro,line cap=round] (0,0) -- (2.6mm,0);}},
}
\newcommand{\simbuno}{\tikz[baseline=-0.6ex]{\draw[relacion](0,0)--(7mm,0);\pic at (4.4mm,0){uno};}}
\newcommand{\simbmuchos}{\tikz[baseline=-0.6ex]{\draw[relacion](0,0)--(7mm,0);\pic at (4.4mm,0){muchos};}}

\usepackage[normalem]{ulem}

\title{Tipos de JOIN}
\subtitle{Unir tablas paso a paso, con ejemplos sencillos}
\date{Semana 09 - Sesión 17}
\author{Rel Guzman}

% Espacio en blanco para completar
\newcommand{\bl}[1][1.9cm]{\underline{\hspace{#1}}}
% Nombre de tabla, siempre en naranja
\newcommand{\tbl}[1]{\textcolor{utpNaranja}{\texttt{#1}}}
% Palabra reservada de SQL dentro de un bloque de código
\newcommand{\kw}[1]{\textcolor{utpRojo}{\textbf{#1}}}
% Parte nueva de un bloque de codigo que crece paso a paso
\newcommand{\nuevo}[1]{{%
  \setlength{\fboxsep}{1.2pt}%
  \begingroup
  \renewcommand{\kw}[1]{\textcolor{utpNegro}{\textbf{##1}}}%
  \colorbox{yellow!45}{\textcolor{utpNegro}{#1}}%
  \endgroup}}
% Comentario dentro de un bloque de código SQL
\newcommand{\com}[1]{\textcolor{utpGris}{-- #1}}
% Marcas de clave primaria y clave foránea
\newcommand{\pk}{{\scriptsize\color{utpRojo}\textbf{PK}}}
\newcommand{\fk}{{\scriptsize\color{utpNaranja}\textbf{FK}}}
% Insignia de práctica
\newcommand{\ins}[1]{\tikz[baseline=-0.5ex]{\node[rounded corners=2pt,fill=utpRojo,text=white,
  font=\bfseries\scriptsize,inner xsep=5pt,inner ysep=2.5pt] {#1};}}
% Etiqueta oscura: nombre del gestor o de la forma de escribir algo
\newcommand{\etiqueta}[1]{\tikz[baseline=-0.5ex]{\node[rounded corners=2pt,fill=utpNegro,text=white,
  font=\bfseries\scriptsize,inner xsep=6pt,inner ysep=3pt] {#1};}}
% Celda de una fila que se quedó sin pareja
\newcommand{\mal}[1]{{\setlength{\fboxsep}{1.2pt}\colorbox{utpRojo!18}{#1}}}
% Valor desconocido
\newcommand{\nulo}{\textcolor{utpGrisMedio}{\textit{NULL}}}
% Bloque de código sobre fondo gris
\newsavebox{\sqlbuf}
\newenvironment{sqlcode}[1][.86\textwidth]
  {\begin{center}\begin{lrbox}{\sqlbuf}\begin{minipage}{#1}\ttfamily\footnotesize\vspace{6pt}\raggedright}
  {\vspace{6pt}\end{minipage}\end{lrbox}\colorbox{utpGrisClaro}{\usebox{\sqlbuf}}\end{center}}
% Sangría dentro de un bloque de código
\newcommand{\ind}{\hspace*{1.6em}}
% Anatomía de una instrucción: cada pieza del código con su pregunta al lado
\tikzset{
  pieza/.style={fill=utpGrisClaro,font=\ttfamily\footnotesize,anchor=west,inner sep=5pt,
    minimum height=0.62cm},
  para/.style={font=\small,anchor=west,text=utpNegro}
}
\newcommand{\piezax}{7.6}
\newcommand{\pieza}[3]{%
  \node[pieza] at (0,#1) {#2};
  \node[para] at (\piezax,#1) {\textcolor{utpRojo}{$\blacktriangleleft$}\ #3};}
% Diapositiva de sección: una palabra bien grande (\seccion[0]{...} la deja sin línea roja)
\newcommand{\seccion}[2][1]{%
\begin{frame}[plain]
  \begin{tikzpicture}[remember picture,overlay]
    \node[anchor=north east,inner sep=0pt]
      at ([xshift=-10mm,yshift=-8mm]current page.north east)
      {\includegraphics[width=43mm]{utp-logo.png}};
    \node[anchor=center,inner sep=0pt] at (current page.center)
      {\resizebox{100mm}{!}{\bfseries\color{utpNegro}#2}};
    \ifnum#1=1
    \node[anchor=north,inner sep=0pt] at ([yshift=-11mm]current page.center)
      {\color{utpRojo}\rule{100mm}{1.6pt}};
    \fi
  \end{tikzpicture}
\end{frame}}
% Diagrama de dos círculos: se pinta la parte que el JOIN conserva.
% \venn[escala]{izquierda}{centro}{derecha}, cada uno con \von o \voff
\newcommand{\von}{utpRojo!45}
\newcommand{\voff}{white}
\newcommand{\venn}[4][1]{%
  \begin{tikzpicture}[scale=#1,baseline=-0.6ex]
    \fill[#2] (0,0) circle (0.62);
    \fill[#4] (0.74,0) circle (0.62);
    \begin{scope}
      \clip (0,0) circle (0.62);
      \fill[#3] (0.74,0) circle (0.62);
    \end{scope}
    \draw[thick,utpNegro] (0,0) circle (0.62);
    \draw[thick,utpNegro] (0.74,0) circle (0.62);
  \end{tikzpicture}}
% Tablas de ejemplo: la tienda
\newcommand{\tablaclientes}{%
  \centering\tbl{clientes}\\[4pt]
  \scriptsize\ttfamily
  \begin{tabular}{rl}
  \toprule
  \normalfont\bfseries id & \normalfont\bfseries nombre \\
  \midrule
  1 & Ana \\
  2 & Luis \\
  3 & María \\
  \bottomrule
  \end{tabular}}
\newcommand{\tablapedidos}{%
  \centering\tbl{pedidos}\\[4pt]
  \scriptsize\ttfamily
  \begin{tabular}{rlr}
  \toprule
  \normalfont\bfseries id & \normalfont\bfseries producto & \normalfont\bfseries cliente\_id \\
  \midrule
  101 & Laptop & 1 \\
  102 & Mouse & 1 \\
  103 & Teclado & 2 \\
  104 & Monitor & \nulo \\
  \bottomrule
  \end{tabular}}
% Tablas de práctica: la universidad
\newcommand{\tablacarreras}{%
  \centering\tbl{carreras}\\[4pt]
  \scriptsize\ttfamily
  \begin{tabular}{rl}
  \toprule
  \normalfont\bfseries id & \normalfont\bfseries nombre \\
  \midrule
  1 & Sistemas \\
  2 & Industrial \\
  3 & Civil \\
  \bottomrule
  \end{tabular}}
\newcommand{\tablaalumnos}{%
  \centering\tbl{alumnos}\\[4pt]
  \scriptsize\ttfamily
  \begin{tabular}{rlr}
  \toprule
  \normalfont\bfseries id & \normalfont\bfseries nombre & \normalfont\bfseries carrera\_id \\
  \midrule
  1 & Paola & 1 \\
  2 & Jorge & 1 \\
  3 & Inés & 3 \\
  4 & Tomás & \nulo \\
  \bottomrule
  \end{tabular}}
% Resultado de un JOIN entre clientes y pedidos: \resultado{filas}
\newcommand{\resultado}[1]{%
  \centering{\scriptsize\color{utpGris}resultado}\\[3pt]
  \footnotesize\ttfamily
  \begin{tabular}{ll}
  \toprule
  \normalfont\bfseries nombre & \normalfont\bfseries producto \\
  \midrule
  #1
  \bottomrule
  \end{tabular}}
\newcommand{\filasinner}{Ana & Laptop \\ Ana & Mouse \\ Luis & Teclado \\}

\begin{document}

% 1 ---------------------------------------------------------------------
{\setbeamertemplate{footline}{}\maketitle}

% 2 ---------------------------------------------------------------------
\begin{frame}[c]
\begin{block}{Logro de la sesión}
Al finalizar, podrás \textbf{unir dos tablas} con \texttt{JOIN} y \textbf{elegir el tipo de JOIN} según lo que quieras conservar.
\end{block}
\end{frame}
% ---------------------------------------------------------------------
\begin{frame}{Lo de hoy: cuatro paradas}
\begin{center}
\begin{tikzpicture}
  \tikzset{parada/.style={cajaverde,font=\footnotesize\bfseries,text width=1.95cm,minimum height=1.15cm}}
  \node[parada] (a) {1. ¿Por qué unir tablas?};
  \node[parada,right=4.5mm of a] (b) {2. \texttt{INNER JOIN}};
  \node[parada,right=4.5mm of b] (c) {3. \texttt{LEFT} y \texttt{RIGHT}};
  \node[parada,right=4.5mm of c] (d) {4. \texttt{FULL}, \texttt{CROSS} y \texttt{SELF}};
  \draw[flecha] (a)--(b); \draw[flecha] (b)--(c); \draw[flecha] (c)--(d);
\end{tikzpicture}
\end{center}
\vspace{0.5cm}
\begin{itemize}\setlength\itemsep{8pt}
  \item Ya sabes consultar \textbf{una tabla}. Hoy consultamos \textbf{dos a la vez}.
  \item Usaremos \textbf{las mismas dos tablas} todo el tiempo: solo cambia el tipo de \texttt{JOIN}.
\end{itemize}
\vspace{0.4cm}
\ojo{Vamos \textbf{de a pocos}: una idea por diapositiva y una práctica después de cada parada.}
\end{frame}
% =======================================================================
\seccion{Unir tablas}
% ---------------------------------------------------------------------
\begin{frame}{Los datos viven en tablas separadas}
\begin{columns}[c,onlytextwidth]
\column{.22\textwidth}
\tablaclientes
\column{.36\textwidth}
\tablapedidos
\column{.38\textwidth}
\small \textbf{Pregunta:} ¿quién pidió el \textbf{Teclado}?
\begin{enumerate}\setlength\itemsep{4pt}
  \item En \tbl{pedidos}: el Teclado tiene \texttt{cliente\_id} = \textbf{2}.
  \item En \tbl{clientes}: el \texttt{id} \textbf{2} es \textbf{Luis}.
\end{enumerate}
\end{columns}
\vspace{0.35cm}
\concepto{La idea}{\small Tú lo resolviste \textbf{mirando dos tablas} y juntando lo que coincide. Eso mismo hace un \texttt{JOIN}, en \textbf{una sola consulta}.}
\vspace{0.1cm}
\analogia{Tienes una hoja con \textbf{códigos y notas}, y otra con \textbf{códigos y nombres}. Las pones lado a lado para saber de quién es cada nota.}
\end{frame}
% ---------------------------------------------------------------------
\begin{frame}{Un \texttt{JOIN}, pieza por pieza}
\begin{center}
\begin{tikzpicture}[y=0.82cm]
  \renewcommand{\piezax}{6.0}
  \pieza{3}{\kw{SELECT} c.nombre, p.producto}{qué columnas quiero, y \textbf{de qué tabla}}
  \pieza{2}{\kw{FROM} clientes c}{la tabla de la \textbf{izquierda}: se escribe primero}
  \pieza{1}{\nuevo{\kw{JOIN}} pedidos p}{la tabla de la \textbf{derecha}: se escribe después}
  \pieza{0}{\nuevo{\kw{ON}} p.cliente\_id = c.id;}{cómo se \textbf{emparejan} las filas}
\end{tikzpicture}
\end{center}
\vspace{0.3cm}
\begin{itemize}\setlength\itemsep{7pt}
  \item \texttt{c} y \texttt{p} son \textbf{apodos} (alias) de las tablas: ahorran escribir el nombre completo.
  \item \texttt{c.nombre} se lee «la columna \texttt{nombre} de \tbl{clientes}».
  \item El \texttt{ON} casi siempre compara la \textbf{clave foránea} con la \textbf{clave primaria}.
\end{itemize}
\end{frame}
% ---------------------------------------------------------------------
\begin{frame}{Al emparejar, algunos se quedan sin pareja}
\begin{center}
\begin{tikzpicture}[y=0.74cm]
  \tikzset{
    fila/.style={draw=utpNegro,thick,fill=utpGrisClaro,font=\footnotesize\ttfamily,
      minimum height=0.52cm,text width=2.7cm,align=left,inner sep=3pt},
    sola/.style={fila,draw=utpRojo!75,fill=utpRojo!12}
  }
  \node[fila] (a) at (0,-0.5) {1 \ Ana};
  \node[fila] (l) at (0,-1.5) {2 \ Luis};
  \node[sola] (m) at (0,-2.5) {3 \ María};
  \node[fila] (p1) at (6.6,0) {Laptop \hfill 1};
  \node[fila] (p2) at (6.6,-1) {Mouse \hfill 1};
  \node[fila] (p3) at (6.6,-2) {Teclado \hfill 2};
  \node[sola] (p4) at (6.6,-3) {Monitor \hfill \nulo};
  \draw[very thick,utpNegro] (a.east)--(p1.west);
  \draw[very thick,utpNegro] (a.east)--(p2.west);
  \draw[very thick,utpNegro] (l.east)--(p3.west);
  \node[font=\small] at (0,0.75) {\tbl{clientes}};
  \node[font=\small] at (6.6,0.75) {\tbl{pedidos}};
  \node[font=\scriptsize,text=utpRojo,left=2mm of m] {no pidió nada};
  \node[font=\scriptsize,text=utpRojo,right=2mm of p4] {no tiene cliente};
\end{tikzpicture}
\end{center}
\vspace{0.1cm}
\begin{itemize}\setlength\itemsep{5pt}
  \item \textbf{Con pareja:} Ana (dos pedidos) y Luis (un pedido).
  \item \textbf{Sin pareja:} \mal{María} y el \mal{Monitor}.
\end{itemize}
\vspace{0.2cm}
\concepto{La única pregunta de hoy}{\small Cada tipo de \texttt{JOIN} responde lo mismo: \textbf{¿qué hago con los que se quedaron sin pareja?}}
\end{frame}
% =======================================================================
\seccion[0]{INNER JOIN}
% ---------------------------------------------------------------------
\begin{frame}{\texttt{INNER JOIN}: solo los que tienen pareja}
\begin{columns}[c,onlytextwidth]
\column{.60\textwidth}
\begin{sqlcode}[\textwidth]
\scriptsize
\kw{SELECT} c.nombre, p.producto\\
\kw{FROM} clientes c\\
\nuevo{\kw{INNER JOIN}} pedidos p \kw{ON} p.cliente\_id = c.id;
\end{sqlcode}
\vspace{-0.15cm}
\centering\venn{\voff}{\von}{\voff}\\[3pt]
{\scriptsize\color{utpGris}solo lo que está en \textbf{las dos} tablas}
\column{.34\textwidth}
\resultado{\filasinner}
\end{columns}
\vspace{0.35cm}
\begin{itemize}\setlength\itemsep{6pt}
  \item Los que no tienen pareja \textbf{quedan fuera}: no aparecen ni María ni el Monitor.
  \item Ana aparece \textbf{dos veces}: una fila por cada pedido suyo.
  \item Escribir solo \texttt{JOIN} es lo mismo que \texttt{INNER JOIN}. Es el más usado.
\end{itemize}
\vspace{0.15cm}
\analogia{Es una \textbf{fiesta solo para parejas}: el que llega solo, no entra.}
\end{frame}
% ---------------------------------------------------------------------
\begin{frame}{\ins{Práctica 1} \texttt{INNER JOIN} en la universidad}
\begin{columns}[c,onlytextwidth]
\column{.24\textwidth}
\tablacarreras
\column{.34\textwidth}
\tablaalumnos
\column{.40\textwidth}
\begin{sqlcode}[\textwidth]
\scriptsize
\kw{SELECT} c.nombre, a.nombre\\
\kw{FROM} carreras c\\
\kw{INNER JOIN} alumnos a\\
\ind \kw{ON} a.carrera\_id = c.id;
\end{sqlcode}
\end{columns}
\vspace{0.2cm}
\begin{center}
\small
\begin{tabular}{>{\raggedright\arraybackslash}p{8.6cm}>{\raggedright\arraybackslash}p{3.4cm}}
\toprule
\bfseries Pregunta & \bfseries Respuesta \\
\midrule
¿Cuántas filas tiene el resultado? & \bl[3.2cm] \\[9pt]
¿Qué carrera \textbf{no} aparece? ¿Por qué? & \bl[3.2cm] \\[9pt]
¿Qué alumno \textbf{no} aparece? ¿Por qué? & \bl[3.2cm] \\
\bottomrule
\end{tabular}
\end{center}
\vspace{0.2cm}
\small \textbf{Para escribir:} anota en tu cuaderno las filas del resultado, una por una.
\end{frame}
% =======================================================================
\seccion{LEFT y RIGHT}
% ---------------------------------------------------------------------
\begin{frame}{\texttt{LEFT JOIN}: toda la tabla izquierda}
\begin{columns}[c,onlytextwidth]
\column{.60\textwidth}
\begin{sqlcode}[\textwidth]
\scriptsize
\kw{SELECT} c.nombre, p.producto\\
\kw{FROM} clientes c\\
\nuevo{\kw{LEFT JOIN}} pedidos p \kw{ON} p.cliente\_id = c.id;
\end{sqlcode}
\vspace{-0.15cm}
\centering\venn{\von}{\von}{\voff}\\[3pt]
{\scriptsize\color{utpGris}\textbf{toda} la izquierda, tenga pareja o no}
\column{.34\textwidth}
\resultado{\filasinner \mal{María} & \mal{\nulo} \\}
\end{columns}
\vspace{0.35cm}
\begin{itemize}\setlength\itemsep{6pt}
  \item La \textbf{izquierda} es la tabla del \texttt{FROM}: \tbl{clientes}. \textbf{Ninguno} de sus clientes se pierde.
  \item María no tiene pedido: su \texttt{producto} se rellena con \nulo.
  \item El Monitor sigue fuera: está en la tabla de la \textbf{derecha}.
\end{itemize}
\vspace{0.15cm}
\analogia{Es la \textbf{lista de asistencia}: salen todos los alumnos, y al que faltó se le deja el casillero vacío.}
\end{frame}
% ---------------------------------------------------------------------
\begin{frame}{\texttt{RIGHT JOIN}: toda la tabla derecha}
\begin{columns}[c,onlytextwidth]
\column{.60\textwidth}
\begin{sqlcode}[\textwidth]
\scriptsize
\kw{SELECT} c.nombre, p.producto\\
\kw{FROM} clientes c\\
\nuevo{\kw{RIGHT JOIN}} pedidos p \kw{ON} p.cliente\_id = c.id;
\end{sqlcode}
\vspace{-0.15cm}
\centering\venn{\voff}{\von}{\von}\\[3pt]
{\scriptsize\color{utpGris}\textbf{toda} la derecha, tenga pareja o no}
\column{.34\textwidth}
\resultado{\filasinner \mal{\nulo} & \mal{Monitor} \\}
\end{columns}
\vspace{0.35cm}
\begin{itemize}\setlength\itemsep{6pt}
  \item La \textbf{derecha} es la tabla del \texttt{JOIN}: \tbl{pedidos}. \textbf{Ningún} pedido se pierde.
  \item El Monitor no tiene cliente: su \texttt{nombre} se rellena con \nulo.
  \item Ahora la que queda fuera es María.
\end{itemize}
\vspace{0.15cm}
\small Es el \textbf{mismo} \texttt{LEFT JOIN}, visto desde el otro lado.
\end{frame}
% ---------------------------------------------------------------------
\begin{frame}{\texttt{LEFT} y \texttt{RIGHT}: lo mismo, al revés}
\begin{columns}[T,onlytextwidth]
\column{.48\textwidth}
\centering\etiqueta{Con RIGHT}
\begin{sqlcode}[\textwidth]
\scriptsize
\kw{FROM} clientes c\\
\kw{RIGHT JOIN} \nuevo{pedidos p}\\
\ind \kw{ON} p.cliente\_id = c.id;
\end{sqlcode}
\column{.48\textwidth}
\centering\etiqueta{Con LEFT}
\begin{sqlcode}[\textwidth]
\scriptsize
\kw{FROM} \nuevo{pedidos p}\\
\kw{LEFT JOIN} clientes c\\
\ind \kw{ON} p.cliente\_id = c.id;
\end{sqlcode}
\end{columns}
\vspace{0.1cm}
\begin{center}
\small Las dos conservan \textbf{todos los pedidos}. Solo cambió \textbf{el orden} en que se escriben las tablas.
\end{center}
\vspace{0.2cm}
\begin{itemize}\setlength\itemsep{7pt}
  \item \textbf{Truco:} pregúntate «¿de qué tabla \textbf{no quiero perder} ninguna fila?».
  \item Pon esa tabla en el \texttt{FROM} y usa \texttt{LEFT JOIN}.
\end{itemize}
\vspace{0.3cm}
\enlarealidad{Casi todos escriben siempre \texttt{LEFT JOIN} y acomodan el orden de las tablas. El \texttt{RIGHT JOIN} se ve muy poco.}
\end{frame}
% ---------------------------------------------------------------------
\begin{frame}{Un clásico: encontrar a los que no tienen pareja}
\small \textbf{Pregunta:} ¿qué clientes \textbf{nunca} han hecho un pedido?
\vspace{0.1cm}
\begin{columns}[c,onlytextwidth]
\column{.60\textwidth}
\begin{sqlcode}[\textwidth]
\scriptsize
\kw{SELECT} c.nombre\\
\kw{FROM} clientes c\\
\kw{LEFT JOIN} pedidos p \kw{ON} p.cliente\_id = c.id\\
\nuevo{\kw{WHERE} p.id \kw{IS NULL};}
\end{sqlcode}
\column{.34\textwidth}
\centering{\scriptsize\color{utpGris}resultado}\\[3pt]
\footnotesize\ttfamily
\begin{tabular}{l}
\toprule
\normalfont\bfseries nombre \\
\midrule
María \\
\bottomrule
\end{tabular}
\end{columns}
\vspace{0.2cm}
\begin{enumerate}\setlength\itemsep{6pt}
  \item El \texttt{LEFT JOIN} trae a \textbf{todos} los clientes; a María le pone \nulo{} en el pedido.
  \item El \texttt{WHERE} se queda \textbf{solo con las filas que tienen ese \nulo}.
\end{enumerate}
\vspace{0.3cm}
\enlarealidad{Así se buscan los alumnos \textbf{sin matrícula}, los productos \textbf{que nadie compró} o los empleados \textbf{sin área}.}
\end{frame}
% ---------------------------------------------------------------------
\begin{frame}{\ins{Práctica 2} \texttt{LEFT} y \texttt{RIGHT} en la universidad}
\begin{columns}[c,onlytextwidth]
\column{.24\textwidth}
\tablacarreras
\column{.34\textwidth}
\tablaalumnos
\column{.40\textwidth}
\begin{sqlcode}[\textwidth]
\scriptsize
\kw{SELECT} c.nombre, a.nombre\\
\kw{FROM} carreras c\\
\bl[1.3cm] \kw{JOIN} alumnos a\\
\ind \kw{ON} a.carrera\_id = c.id;
\end{sqlcode}
\end{columns}
\vspace{0.2cm}
\begin{center}
\small
\renewcommand{\arraystretch}{1.1}
\begin{tabular}{>{\raggedright\arraybackslash}p{4.4cm}>{\raggedright\arraybackslash}p{3.6cm}>{\raggedright\arraybackslash}p{3.6cm}}
\toprule
& \bfseries Con \texttt{LEFT} & \bfseries Con \texttt{RIGHT} \\
\midrule
¿Cuántas filas salen? & \bl[3.0cm] & \bl[3.0cm] \\[9pt]
¿Qué fila lleva un \nulo? & \bl[3.0cm] & \bl[3.0cm] \\[9pt]
¿Quién queda fuera? & \bl[3.0cm] & \bl[3.0cm] \\
\bottomrule
\end{tabular}
\end{center}
\vspace{0.2cm}
\small \textbf{Para escribir:} la consulta que muestra las carreras que \textbf{no tienen alumnos}.
\end{frame}
% =======================================================================
\seccion{FULL, CROSS y SELF}
% ---------------------------------------------------------------------
\begin{frame}{\texttt{FULL OUTER JOIN}: nadie queda fuera}
\begin{columns}[c,onlytextwidth]
\column{.60\textwidth}
\begin{sqlcode}[\textwidth]
\scriptsize
\kw{SELECT} c.nombre, p.producto\\
\kw{FROM} clientes c\\
\nuevo{\kw{FULL OUTER JOIN}} pedidos p \kw{ON} p.cliente\_id = c.id;
\end{sqlcode}
\vspace{-0.15cm}
\centering\venn{\von}{\von}{\von}\\[3pt]
{\scriptsize\color{utpGris}\textbf{todo} lo de las dos tablas}
\column{.34\textwidth}
\resultado{\filasinner \mal{María} & \mal{\nulo} \\ \mal{\nulo} & \mal{Monitor} \\}
\end{columns}
\vspace{0.35cm}
\begin{itemize}\setlength\itemsep{6pt}
  \item Es un \texttt{LEFT} y un \texttt{RIGHT} \textbf{juntos}: aparecen María \textbf{y} el Monitor.
  \item A cada uno se le rellena con \nulo{} el lado que le falta.
\end{itemize}
\vspace{0.2cm}
\ojo{\textbf{MySQL no tiene} \texttt{FULL OUTER JOIN}. Ahí se arma juntando un \texttt{LEFT JOIN} y un \texttt{RIGHT JOIN} con \texttt{UNION}.}
\end{frame}
% ---------------------------------------------------------------------
\begin{frame}{\texttt{CROSS JOIN}: todos con todos}
\begin{columns}[c,onlytextwidth]
\column{.16\textwidth}
\centering\tbl{tallas}\\[4pt]
\scriptsize\ttfamily
\begin{tabular}{l}
\toprule
\normalfont\bfseries talla \\
\midrule
S \\ M \\
\bottomrule
\end{tabular}
\column{.18\textwidth}
\centering\tbl{colores}\\[4pt]
\scriptsize\ttfamily
\begin{tabular}{l}
\toprule
\normalfont\bfseries color \\
\midrule
rojo \\ azul \\ negro \\
\bottomrule
\end{tabular}
\column{.38\textwidth}
\begin{sqlcode}[\textwidth]
\scriptsize
\kw{SELECT} t.talla, c.color\\
\kw{FROM} tallas t\\
\nuevo{\kw{CROSS JOIN}} colores c;
\end{sqlcode}
\column{.24\textwidth}
\centering\scriptsize\ttfamily
\begin{tabular}{ll}
\toprule
\normalfont\bfseries talla & \normalfont\bfseries color \\
\midrule
S & rojo \\ S & azul \\ S & negro \\
M & rojo \\ M & azul \\ M & negro \\
\bottomrule
\end{tabular}
\end{columns}
\vspace{0.3cm}
\begin{itemize}\setlength\itemsep{6pt}
  \item \textbf{No lleva} \texttt{ON}: no empareja nada, combina \textbf{cada fila con todas} las de la otra tabla.
  \item Las filas se \textbf{multiplican}: 2 tallas $\times$ 3 colores = \textbf{6 filas}.
\end{itemize}
\vspace{0.2cm}
\ojo{Con tablas grandes explota: 1\,000 $\times$ 1\,000 = \textbf{1\,000\,000} de filas. Se usa poco, y a propósito.}
\end{frame}
% ---------------------------------------------------------------------
\begin{frame}{\texttt{SELF JOIN}: una tabla unida consigo misma}
\begin{columns}[c,onlytextwidth]
\column{.27\textwidth}
\centering\tbl{empleados}\\[4pt]
\scriptsize\ttfamily
\begin{tabular}{rlr}
\toprule
\normalfont\bfseries id & \normalfont\bfseries nombre & \normalfont\bfseries jefe\_id \\
\midrule
1 & Rosa & \nulo \\
2 & Carlos & 1 \\
3 & Elena & 1 \\
4 & Diego & 2 \\
\bottomrule
\end{tabular}
\column{.44\textwidth}
\begin{sqlcode}[\textwidth]
\scriptsize
\kw{SELECT} e.nombre \kw{AS} empleado,\\
\ind\ind\ \ j.nombre \kw{AS} jefe\\
\kw{FROM} empleados \nuevo{e}\\
\kw{JOIN} empleados \nuevo{j} \kw{ON} e.jefe\_id = j.id;
\end{sqlcode}
\column{.25\textwidth}
\centering\scriptsize\ttfamily
\begin{tabular}{ll}
\toprule
\normalfont\bfseries empleado & \normalfont\bfseries jefe \\
\midrule
Carlos & Rosa \\
Elena & Rosa \\
Diego & Carlos \\
\bottomrule
\end{tabular}
\end{columns}
\vspace{0.3cm}
\begin{itemize}\setlength\itemsep{6pt}
  \item El jefe \textbf{también es un empleado}: está en la misma tabla.
  \item \textbf{No es una palabra nueva}: es un \texttt{JOIN} de siempre, con la tabla escrita \textbf{dos veces}.
  \item Los apodos son \textbf{obligatorios}: \texttt{e} es el empleado y \texttt{j} es su jefe.
\end{itemize}
\vspace{0.15cm}
\small Rosa no aparece como empleado porque no tiene jefe. ¿Con qué \texttt{JOIN} sí aparecería?
\end{frame}
% ---------------------------------------------------------------------
\begin{frame}{Todos los \texttt{JOIN} en una tabla}
\begin{center}
\renewcommand{\arraystretch}{1.75}
\small
\begin{tabular}{>{\centering\arraybackslash}m{1.7cm}>{\ttfamily\bfseries\raggedright\arraybackslash}m{3.0cm}>{\raggedright\arraybackslash}m{6.1cm}>{\centering\arraybackslash}m{1.5cm}}
\toprule
& \normalfont\bfseries Tipo & \bfseries Qué conserva & \bfseries Filas \\
\midrule
\venn[0.42]{\voff}{\von}{\voff} & INNER JOIN & solo los que tienen pareja & 3 \\
\venn[0.42]{\von}{\von}{\voff} & LEFT JOIN & toda la izquierda & 4 \\
\venn[0.42]{\voff}{\von}{\von} & RIGHT JOIN & toda la derecha & 4 \\
\venn[0.42]{\von}{\von}{\von} & FULL OUTER JOIN & todo, de los dos lados & 5 \\
$\times$ & CROSS JOIN & todas las combinaciones, sin \texttt{ON} & 12 \\
\bottomrule
\end{tabular}
\end{center}
\vspace{0.1cm}
\small Las filas son las de \tbl{clientes} (3) con \tbl{pedidos} (4). El \texttt{SELF JOIN} no está: es cualquiera de estos, con la misma tabla dos veces.
\end{frame}
% ---------------------------------------------------------------------
\begin{frame}{\ins{Práctica 3} ¿Qué \texttt{JOIN} usarías?}
\begin{center}
\small
\renewcommand{\arraystretch}{1.15}
\begin{tabular}{>{\raggedright\arraybackslash}p{9.6cm}>{\raggedright\arraybackslash}p{2.8cm}}
\toprule
\bfseries Lo que piden & \bfseries \texttt{JOIN} \\
\midrule
Cada alumno con su carrera; \textbf{solo} los que tienen carrera. & \bl[2.6cm] \\[8pt]
\textbf{Todas} las carreras, tengan alumnos o no. & \bl[2.6cm] \\[8pt]
Carreras y alumnos, \textbf{sin dejar a nadie} fuera. & \bl[2.6cm] \\[8pt]
Todas las combinaciones de 3 polos y 4 colores. & \bl[2.6cm] \\[8pt]
Cada empleado con el nombre de \textbf{su jefe}. & \bl[2.6cm] \\
\bottomrule
\end{tabular}
\end{center}
\vspace{0.3cm}
\begin{itemize}\setlength\itemsep{5pt}
  \item ¿Cuántas filas salen en la cuarta? \bl[1.5cm]
  \item Con \tbl{carreras} (3) y \tbl{alumnos} (4), ¿cuántas filas da el \texttt{FULL OUTER JOIN}? \bl[1.5cm]
\end{itemize}
\end{frame}
% ---------------------------------------------------------------------
\begin{frame}[plain]
  \begin{tikzpicture}[remember picture,overlay]
    \node[anchor=north west,font=\bfseries\fontsize{24}{28}\selectfont]
      at ([xshift=12mm,yshift=-9mm]current page.north west) {Preguntas};
    \node[anchor=north east,inner sep=0pt]
      at ([xshift=-10mm,yshift=-8mm]current page.north east)
      {\includegraphics[width=43mm]{utp-logo.png}};
    \node[anchor=center,inner sep=0pt]
      at ([yshift=-5mm]current page.center)
      {\includegraphics[width=78mm]{preguntas-ilustracion.png}};
  \end{tikzpicture}
\end{frame}
% ---------------------------------------------------------------------
\begin{frame}{Cierre: lo de hoy en una tabla}
\begin{center}
\renewcommand{\arraystretch}{1.45}
\begin{tabular}{>{\raggedright\bfseries}p{3.4cm}>{\raggedright\arraybackslash}p{9.0cm}}
\toprule
& \bfseries En una frase \\
\midrule
\texttt{JOIN} & Junta filas de dos tablas que coinciden en el \texttt{ON}. \\
\texttt{INNER} & Solo los que tienen pareja. \\
\texttt{LEFT} / \texttt{RIGHT} & Toda una tabla; lo que falta del otro lado queda en \nulo. \\
\texttt{FULL} / \texttt{CROSS} / \texttt{SELF} & Todo de ambas; todos con todos; la tabla consigo misma. \\
\bottomrule
\end{tabular}
\end{center}
\vspace{0.3cm}
\small La pregunta para elegir un \texttt{JOIN}: \textbf{¿qué hago con los que no tienen pareja?}
\end{frame}
% ---------------------------------------------------------------------
\begin{frame}{Referencias}
\footnotesize
Coronel, C., Morris, S., \& Rob, P. (2011). \textit{Bases de datos: diseño, implementación y administración} (9.ª ed.). Cengage Learning.

\vspace{0.3cm}
Elmasri, R., \& Navathe, S. (2007). \textit{Fundamentos de sistemas de bases de datos} (5.ª ed.). Addison Wesley.

\vspace{0.3cm}
Silberschatz, A., Korth, H. F., \& Sudarshan, S. (2020). \textit{Database system concepts} (7th ed.). McGraw-Hill.

\vspace{0.3cm}
Oracle. (s.\,f.). \textit{MySQL reference manual: JOIN clause}. \texttt{https://dev.mysql.com/doc/refman/8.4/en/join.html}
\end{frame}
% ---------------------------------------------------------------------
% Cierre UTP
{
  \setbeamertemplate{background canvas}{%
    \includegraphics[width=\paperwidth,height=\paperheight]{gracias-final.pdf}%
  }
  \setbeamertemplate{footline}{}
  \begin{frame}[plain]
  \end{frame}
}

\end{document}
